\section{Distributional robustness and cross-validation}
\label{sec:distributional_robustness}

\subsection{Motivation and claim boundary}

The headline design fixes a train/test split and asks whether the option-only optimizer,
cost layers, and settlement gates survive a single frozen out-of-sample path. This section
asks a different question: how fragile are those conclusions to alternative month
orderings, refit samples, and synthetic repricing worlds? The answer is reported through
blocked cross-validation, combinatorial purged cross-validation (CPCV), resampled
universes, and repriced Monte Carlo universes. The tables carry the numerical evidence;
the text here records the design and the claim boundary.

\begin{center}
\fbox{\begin{minipage}{0.88\textwidth}
\textbf{\color{journalgreen} Honesty box}\par
CPCV is deliberately non-PIT. By construction it trains some folds on observations after
the fold being tested, so it is a distributional-robustness and overfitting diagnostic
(path distribution and PBO), not a tradable out-of-sample claim. The headline claim
boundary from the results and conclusion sections is unchanged.\par
\vspace{2pt}
The contract-universe eligibility filter is also fixed at the headline anchor: contracts
must have at least 36 training observations through the global
\texttt{TRAIN\_END}. That same eligible universe is reused in every fold. This is benign
for the diagnostic because the fold exercise is not a tradable signal, but it is disclosed
because the fold membership is not re-estimated from each fold's own training set.\par
\vspace{2pt}
Per-fold CV refits cover the six option strategies that the option-contract evaluator can
score: Equity-option Greek Markowitz, Greek Markowitz + VIX, Beta/delta-neutral + VIX,
Equal premium, Equal risk, and the VIX hedge sleeve. Cost-aware Sortino + VIX is excluded
from per-fold CV because its per-fold cost-aware solver is expensive; it is included in
the fixed-weight Monte Carlo layer. Delta-matched equities and Underlying Markowitz remain
headline-only in CV because their weight vectors are underlying weights, not option
contract weights.
\end{minipage}}
\end{center}

\subsection{Blocked folds and CPCV design}

The realized monthly return grid is split into 12 contiguous groups. The blocked k-fold
diagnostic tests one group at a time. The CPCV diagnostic tests \(k=2\) groups at a time,
giving \(\binom{12}{2}=66\) splits. Complete CPCV paths are then assembled so that each
path includes each return month exactly once; with this design the complete-path count is
\(\binom{11}{1}=11\), while incomplete paths are kept in the ledger but excluded from the
path distribution.

The purge rule is intentionally simple and auditable. Option labels' payoffs realize at
most 44 calendar days after the decision date. Each fold uses purge \(=1\) month and
embargo \(=1\) month, which creates at least a two-snapshot train/test gap. In the shipped
monthly grid that is at least 57 calendar days, and \(57>44\), so the payoff window of a
test label cannot overlap a training label kept by the fold. The verifier recomputes this
schedule invariant from the fold ledger rather than trusting this paragraph.

Within each fold the option strategies are refit on the fold's training months and scored
on the fold's test months. Strategy metrics use limited-liability absorption at zero, so a
path that loses all wealth remains defaulted rather than sign-flipping. Table~\ref{tab:cv_fold_performance}
reports the blocked k-fold refits, and Table~\ref{tab:cv_cpcv_distribution} reports the
CPCV path distribution against the headline realized path.

\clearpage

\begin{table}[H]
\centering
\scriptsize
\IfFileExists{tables/cv_fold_performance.tex}
  {\resizebox{\textwidth}{!}{\begin{tabular}{lllllll}
\toprule
Fold & Equity-option Greek Markowitz & Greek Markowitz + VIX & Beta/delta-neutral + VIX & Equal premium & Equal risk & VIX hedge sleeve \\
\midrule
kfold\_00 & -0.45 / -1.69 & -0.78 / -2.21 & -0.47 / -2.29 & -0.13 / -0.91 & 0.23 / -0.99 & -6.64 / -2.38 \\
kfold\_01 & 0.11 / -0.51 & -2.53 / -1.02 & -2.49 / -1.19 & 0.47 / 0.75 & 0.34 / 0.82 & -0.60 / -1.56 \\
kfold\_02 & 3.37 / 2.04 & 2.42 / 1.89 & 2.07 / 1.44 & -0.91 / 0.07 & 0.12 / 0.68 & -2.27 / 0.00 \\
kfold\_03 & -0.67 / -0.73 & 0.06 / -2.36 & -1.10 / -2.72 & 0.35 / -0.06 & 0.62 / 0.12 & 0.04 / 0.36 \\
kfold\_04 & 0.90 / 0.69 & -0.17 / -0.09 & 0.58 / 0.41 & 0.20 / -0.49 & 0.42 / -0.41 & -1.03 / 0.23 \\
kfold\_05 & 2.66 / 2.16 & 1.42 / 0.13 & 0.91 / -0.26 & 2.15 / 1.93 & 2.69 / 2.21 & 0.72 / 0.96 \\
kfold\_06 & 0.69 / 0.56 & 1.24 / 0.61 & 1.06 / 0.16 & -1.27 / -3.01 & -1.49 / -2.58 & -5.82 / -1.38 \\
kfold\_07 & -0.73 / -1.49 & -0.05 / -0.40 & 1.13 / 0.75 & -0.16 / -0.85 & -0.36 / -0.91 & -4.70 / -1.63 \\
kfold\_08 & 1.53 / 2.37 & 1.25 / 1.70 & 1.60 / 2.70 & 0.96 / 0.22 & 1.19 / -0.05 & -20.28 / -2.50 \\
kfold\_09 & 2.49 / 2.59 & 3.78 / 1.42 & 2.86 / 1.42 & 0.65 / 1.21 & 1.30 / 1.34 & -19.89 / -1.42 \\
kfold\_10 & 1.53 / 1.29 & 2.00 / 1.07 & 2.29 / 1.24 & -0.12 / -0.04 & 0.03 / -0.64 & -4.36 / -2.71 \\
kfold\_11 & -0.19 / -0.19 & 0.34 / 0.13 & 0.67 / -0.05 & -0.56 / -0.02 & -0.62 / -0.22 & -5.40 / -1.85 \\
\bottomrule
\end{tabular}
}}
  {\begin{tabular}{llll}\toprule Fold & Strategy & Gross Sharpe & Status \\\midrule Missing & Run \texttt{make robustness} & -- & -- \\\bottomrule\end{tabular}}
\caption{Blocked k-fold refit performance. Rows are diagnostics for refitted option
strategies and do not alter the headline out-of-sample claim.}
\label{tab:cv_fold_performance}
\end{table}

\begin{figure}[H]
\centering
\IfFileExists{figures/cv_fold_sharpe_heatmap.pdf}
  {\includegraphics[width=0.86\textwidth]{cv_fold_sharpe_heatmap.pdf}}
  {\fbox{\parbox{0.78\textwidth}{Run \texttt{python -m research.papers.option_only_markowitz.analysis.run_empirics --stage robustness-figures} to generate the fold Sharpe heatmap.}}}
\caption{Blocked k-fold gross Sharpe by strategy and chronological test window. The
color scale is clipped at $\pm 3$ so cross-fold variation in the equity-option
strategies stays visible; cell annotations report unclipped values (the hedge sleeve
reaches $-20$ in the 2022--2024 folds).}
\label{fig:cv_fold_sharpe_heatmap}
\end{figure}

\clearpage

\begin{table}[H]
\centering
\scriptsize
\IfFileExists{tables/cv_cpcv_distribution.tex}
  {\resizebox{\textwidth}{!}{\begin{tabular}{lrrrrrrrrrr}
\toprule
Strategy & Realized Gross Sharpe & CPCV Gross P05 & CPCV Gross P50 & CPCV Gross P95 & Realized Full Spread Sharpe & CPCV Full Spread P05 & CPCV Full Spread P50 & CPCV Full Spread P95 & PBO Gross & PBO Full Spread \\
\midrule
Beta/delta-neutral + VIX & 1.414 & 0.587 & 0.710 & 0.887 & 0.678 & -0.032 & 0.046 & 0.165 & 0.192 & 0.205 \\
Cost-aware Sortino + VIX & 1.014 &  &  &  & 0.047 &  &  &  & 0.192 & 0.205 \\
Delta-matched equities & 1.176 &  &  &  & 1.176 &  &  &  & 0.192 & 0.205 \\
Equal premium & -0.030 & 0.236 & 0.236 & 0.236 & -0.156 & 0.156 & 0.156 & 0.156 & 0.192 & 0.205 \\
Equal risk & 0.082 & 0.410 & 0.465 & 0.499 & -0.107 & 0.156 & 0.195 & 0.274 & 0.192 & 0.205 \\
Equity-option Greek Markowitz & 0.842 & 0.788 & 0.936 & 1.009 & 0.913 & 0.545 & 0.703 & 0.782 & 0.192 & 0.205 \\
Greek Markowitz + VIX & 1.374 & 0.602 & 0.766 & 0.954 & 0.294 & -0.052 & 0.074 & 0.164 & 0.192 & 0.205 \\
Underlying Markowitz & 0.909 &  &  &  & 0.909 &  &  &  & 0.192 & 0.205 \\
VIX hedge sleeve & -6.429 & -0.404 & -0.404 & -0.404 & -1.777 & 0.069 & 0.069 & 0.069 & 0.192 & 0.205 \\
\bottomrule
\end{tabular}
}}
  {\begin{tabular}{llll}\toprule Strategy & Realized Sharpe & CPCV P50 & PBO \\\midrule Missing & Run \texttt{make robustness} & -- & -- \\\bottomrule\end{tabular}}
\caption{CPCV complete-path distribution. The table compares realized-path metrics with
the complete CPCV path distribution under gross and full-spread post-cost bases.}
\label{tab:cv_cpcv_distribution}
\end{table}

\begin{figure}[H]
\centering
\IfFileExists{figures/cv_cpcv_sharpe_distribution.pdf}
  {\includegraphics[width=0.86\textwidth]{cv_cpcv_sharpe_distribution.pdf}}
  {\fbox{\parbox{0.78\textwidth}{Run \texttt{python -m research.papers.option_only_markowitz.analysis.run_empirics --stage robustness-figures} to generate the CPCV Sharpe distribution figure.}}}
\caption{Complete-path CPCV Sharpe distributions with the realized static-split gross
Sharpe marked for each CV strategy. The standalone VIX hedge sleeve (realized gross
Sharpe $-6.4$) is omitted for scale; its fold-level values appear in
Table~\ref{tab:cv_fold_performance} and Figure~\ref{fig:cv_fold_sharpe_heatmap}.}
\label{fig:cv_cpcv_sharpe_distribution}
\end{figure}

\clearpage

\subsection{Probability of backtest overfitting}

The combinatorially symmetric cross-validation (CSCV) probability of backtest overfitting
(PBO) partitions the strategy set into
in-sample and out-of-sample halves, selects the best in-sample strategy, and records the
logit rank it attains out of sample. PBO is the fraction of splits with below-median
out-of-sample rank. In this paper the gross variant selects on gross in-sample Sharpe and
evaluates gross out-of-sample Sharpe. The full-spread variant selects on gross in-sample
Sharpe and evaluates full-spread post-cost out-of-sample Sharpe, so it asks whether
pre-cost selection survives a conservative executable-cost spread layer.

\subsection{Regime-tagged performance}

CPCV path months are also tagged by VIX terciles and by three named event windows:
Volmageddon in February--March 2018, the COVID crash in February--April 2020, and the
2022 tightening bear market. The regime table is descriptive; the labels summarize path
outcomes and do not enter expected returns, constraints, or strategy selection.

\begin{table}[H]
\centering
\scriptsize
\IfFileExists{tables/cv_regime_performance.tex}
  {\resizebox{\textwidth}{!}{\begin{tabular}{llllrrrrr}
\toprule
Regime family & Strategy & Group & Metric & Estimate & CI lo & CI hi & N & Seed \\
\midrule
VIX tercile & Beta/delta-neutral + VIX & High VIX & Sharpe & 0.704 & 0.207 & 1.653 & 43 & 20260625 \\
VIX tercile & Beta/delta-neutral + VIX & Low VIX & Sharpe & 0.625 & -0.706 & 1.781 & 43 & 20260625 \\
VIX tercile & Beta/delta-neutral + VIX & Mid VIX & Sharpe & 0.986 & 0.169 & 1.795 & 43 & 20260625 \\
VIX tercile & Equal premium & High VIX & Sharpe & 0.529 & -0.905 & 1.399 & 43 & 20260625 \\
VIX tercile & Equal premium & Low VIX & Sharpe & 0.365 & -0.263 & 0.898 & 43 & 20260625 \\
VIX tercile & Equal premium & Mid VIX & Sharpe & -0.304 & -1.046 & 0.178 & 43 & 20260625 \\
VIX tercile & Equal risk & High VIX & Sharpe & 0.471 & -0.962 & 1.368 & 43 & 20260625 \\
VIX tercile & Equal risk & Low VIX & Sharpe & 0.843 & 0.319 & 1.398 & 43 & 20260625 \\
VIX tercile & Equal risk & Mid VIX & Sharpe & 0.066 & -0.491 & 0.494 & 43 & 20260625 \\
VIX tercile & Equity-option Greek Markowitz & High VIX & Sharpe & 0.041 & -1.039 & 0.977 & 43 & 20260625 \\
VIX tercile & Equity-option Greek Markowitz & Low VIX & Sharpe & 2.560 & 1.954 & 3.218 & 43 & 20260625 \\
VIX tercile & Equity-option Greek Markowitz & Mid VIX & Sharpe & 0.871 & 0.304 & 1.411 & 43 & 20260625 \\
VIX tercile & Greek Markowitz + VIX & High VIX & Sharpe & 0.418 & -0.260 & 1.202 & 43 & 20260625 \\
VIX tercile & Greek Markowitz + VIX & Low VIX & Sharpe & 1.065 & -0.115 & 2.165 & 43 & 20260625 \\
VIX tercile & Greek Markowitz + VIX & Mid VIX & Sharpe & 1.031 & 0.186 & 1.855 & 43 & 20260625 \\
VIX tercile & VIX hedge sleeve & High VIX & Sharpe & 0.215 & -6.796 & 0.734 & 43 & 20260625 \\
VIX tercile & VIX hedge sleeve & Low VIX & Sharpe & -4.677 & -7.507 & -3.120 & 43 & 20260625 \\
VIX tercile & VIX hedge sleeve & Mid VIX & Sharpe & -5.496 & -8.274 & -4.210 & 43 & 20260625 \\
Event window & Beta/delta-neutral + VIX & 2022 tightening bear & Sharpe & 1.444 & -0.062 & 3.644 & 9 & 20260625 \\
Event window & Beta/delta-neutral + VIX & COVID crash & Sharpe &  &  &  & 3 & 20260625 \\
Event window & Beta/delta-neutral + VIX & Volmageddon Feb 2018 & Sharpe &  &  &  & 1 & 20260625 \\
Event window & Equal premium & 2022 tightening bear & Sharpe & -0.156 & -3.237 & 1.003 & 9 & 20260625 \\
Event window & Equal premium & COVID crash & Sharpe &  &  &  & 3 & 20260625 \\
Event window & Equal premium & Volmageddon Feb 2018 & Sharpe &  &  &  & 1 & 20260625 \\
Event window & Equal risk & 2022 tightening bear & Sharpe & -0.899 & -3.112 & 0.310 & 9 & 20260625 \\
Event window & Equal risk & COVID crash & Sharpe &  &  &  & 3 & 20260625 \\
Event window & Equal risk & Volmageddon Feb 2018 & Sharpe &  &  &  & 1 & 20260625 \\
Event window & Equity-option Greek Markowitz & 2022 tightening bear & Sharpe & -2.574 & -5.433 & -0.978 & 9 & 20260625 \\
Event window & Equity-option Greek Markowitz & COVID crash & Sharpe &  &  &  & 3 & 20260625 \\
Event window & Equity-option Greek Markowitz & Volmageddon Feb 2018 & Sharpe &  &  &  & 1 & 20260625 \\
Event window & Greek Markowitz + VIX & 2022 tightening bear & Sharpe & -0.682 & -1.219 & -0.140 & 9 & 20260625 \\
Event window & Greek Markowitz + VIX & COVID crash & Sharpe &  &  &  & 3 & 20260625 \\
Event window & Greek Markowitz + VIX & Volmageddon Feb 2018 & Sharpe &  &  &  & 1 & 20260625 \\
Event window & VIX hedge sleeve & 2022 tightening bear & Sharpe & -3.963 & -5.637 & -2.927 & 9 & 20260625 \\
Event window & VIX hedge sleeve & COVID crash & Sharpe &  &  &  & 3 & 20260625 \\
Event window & VIX hedge sleeve & Volmageddon Feb 2018 & Sharpe &  &  &  & 1 & 20260625 \\
\bottomrule
\end{tabular}
}}
  {\begin{tabular}{llll}\toprule Regime & Strategy & Metric & Estimate \\\midrule Missing & Run \texttt{make robustness} & -- & -- \\\bottomrule\end{tabular}}
\caption{CPCV regime-tagged performance by VIX tercile and event window.}
\label{tab:cv_regime_performance}
\end{table}

\clearpage

\subsection{Resampled universes}

The first Monte Carlo tier resamples month indices rather than individual strategy values:
one circular block-bootstrap path is applied jointly to every strategy return frame, so
cross-sectional dependence is preserved. The default fixed-weight tier uses 1,000 paths
and evaluates the train-estimated portfolios. This measures robustness of the fixed
portfolio actually used in the headline exercise.

A separate refit tier uses 200 pseudo-training paths for Greek Markowitz + VIX and then
scores each refit on the fixed realized test window. This is a procedure and
estimation-risk diagnostic, not a replacement for the fixed-weight path distribution. In
that refit tier, SPY beta and stress augmentation use slot-calendar dates while the moment
inputs are slot-relabeled resampled months. The approximation can affect constraint bounds
only; it never enters expected returns or realized P\&L.

\begin{table}[H]
\centering
\scriptsize
\IfFileExists{tables/mc_resampled_universes.tex}
  {\resizebox{\textwidth}{!}{\begin{tabular}{lllrrrrrrrrrr}
\toprule
Universe Family & Basis & Strategy & Realized Value & Path P05 Sharpe & Path P25 Sharpe & Path P50 Sharpe & Path P75 Sharpe & Path P95 Sharpe & Path P50 Max Drawdown & P Sharpe Less Than 0 & P Max Drawdown Less Than -0.5 & P Default \\
\midrule
resampled & full\_spread\_post\_cost & Beta/delta-neutral + VIX & 0.678 & -0.004 & 0.401 & 0.658 & 0.919 & 1.305 & -0.513 & 0.051 & 0.535 & 0.000 \\
resampled & full\_spread\_post\_cost & Cost-aware Sortino + VIX & 0.047 & -0.817 & -0.289 & 0.050 & 0.363 & 0.828 & -1.000 & 0.459 & 1.000 & 0.616 \\
resampled & full\_spread\_post\_cost & Delta-matched equities & 1.176 & 0.418 & 0.885 & 1.215 & 1.613 & 2.183 & -0.774 & 0.007 & 0.902 & 0.000 \\
resampled & full\_spread\_post\_cost & Equal premium & -0.156 & -0.796 & -0.375 & -0.158 & 0.038 & 0.291 & -0.984 & 0.703 & 1.000 & 0.000 \\
resampled & full\_spread\_post\_cost & Equal risk & -0.107 & -0.677 & -0.339 & -0.107 & 0.126 & 0.412 & -0.914 & 0.621 & 0.996 & 0.000 \\
resampled & full\_spread\_post\_cost & Equity-option Greek Markowitz & 0.913 & 0.160 & 0.614 & 0.894 & 1.175 & 1.562 & -0.424 & 0.024 & 0.269 & 0.000 \\
resampled & full\_spread\_post\_cost & Greek Markowitz + VIX & 0.294 & -0.399 & 0.018 & 0.269 & 0.522 & 0.866 & -0.655 & 0.236 & 0.858 & 0.000 \\
resampled & full\_spread\_post\_cost & Underlying Markowitz & 0.909 & 0.133 & 0.604 & 0.939 & 1.276 & 1.763 & -0.232 & 0.026 & 0.013 & 0.000 \\
resampled & full\_spread\_post\_cost & VIX hedge sleeve & -1.777 & -2.275 & -2.016 & -1.830 & -1.615 & -1.279 & -1.000 & 1.000 & 1.000 & 0.623 \\
resampled & gross & Beta/delta-neutral + VIX & 1.414 & 0.737 & 1.141 & 1.428 & 1.728 & 2.159 & -0.547 & 0.000 & 0.585 & 0.000 \\
resampled & gross & Cost-aware Sortino + VIX & 1.014 & 0.431 & 0.782 & 1.006 & 1.248 & 1.574 & -1.000 & 0.008 & 1.000 & 0.905 \\
resampled & gross & Delta-matched equities & 1.176 & 0.418 & 0.885 & 1.215 & 1.613 & 2.183 & -0.774 & 0.007 & 0.902 & 0.000 \\
resampled & gross & Equal premium & -0.030 & -0.693 & -0.263 & -0.033 & 0.197 & 0.471 & -1.000 & 0.546 & 1.000 & 0.000 \\
resampled & gross & Equal risk & 0.082 & -0.602 & -0.169 & 0.085 & 0.334 & 0.669 & -0.982 & 0.409 & 1.000 & 0.000 \\
resampled & gross & Equity-option Greek Markowitz & 0.842 & 0.085 & 0.529 & 0.864 & 1.204 & 1.657 & -0.840 & 0.038 & 0.983 & 0.000 \\
resampled & gross & Greek Markowitz + VIX & 1.374 & 0.643 & 1.068 & 1.391 & 1.721 & 2.236 & -0.602 & 0.002 & 0.795 & 0.000 \\
resampled & gross & Underlying Markowitz & 0.909 & 0.133 & 0.604 & 0.939 & 1.276 & 1.763 & -0.232 & 0.026 & 0.013 & 0.000 \\
resampled & gross & VIX hedge sleeve & -6.429 & -9.234 & -7.291 & -6.400 & -5.730 & -4.993 & -1.000 & 1.000 & 1.000 & 1.000 \\
resampled\_stratified & full\_spread\_post\_cost & Beta/delta-neutral + VIX & 0.678 & 0.001 & 0.408 & 0.656 & 0.938 & 1.317 & -0.584 & 0.049 & 0.703 & 0.000 \\
resampled\_stratified & full\_spread\_post\_cost & Cost-aware Sortino + VIX & 0.047 & -0.665 & -0.234 & 0.044 & 0.326 & 0.697 & -1.000 & 0.463 & 1.000 & 0.691 \\
resampled\_stratified & full\_spread\_post\_cost & Delta-matched equities & 1.176 & 0.810 & 1.023 & 1.196 & 1.376 & 1.617 & -0.832 & 0.000 & 0.991 & 0.000 \\
resampled\_stratified & full\_spread\_post\_cost & Equal premium & -0.156 & -0.848 & -0.393 & -0.163 & 0.032 & 0.297 & -0.986 & 0.700 & 1.000 & 0.000 \\
resampled\_stratified & full\_spread\_post\_cost & Equal risk & -0.107 & -0.773 & -0.350 & -0.124 & 0.123 & 0.426 & -0.923 & 0.639 & 0.998 & 0.000 \\
resampled\_stratified & full\_spread\_post\_cost & Equity-option Greek Markowitz & 0.913 & 0.315 & 0.667 & 0.915 & 1.136 & 1.440 & -0.380 & 0.005 & 0.134 & 0.000 \\
resampled\_stratified & full\_spread\_post\_cost & Greek Markowitz + VIX & 0.294 & -0.340 & 0.047 & 0.273 & 0.524 & 0.880 & -0.673 & 0.222 & 0.881 & 0.000 \\
resampled\_stratified & full\_spread\_post\_cost & Underlying Markowitz & 0.909 & 0.472 & 0.748 & 0.929 & 1.089 & 1.325 & -0.260 & 0.001 & 0.000 & 0.000 \\
resampled\_stratified & full\_spread\_post\_cost & VIX hedge sleeve & -1.777 & -2.139 & -1.932 & -1.797 & -1.651 & -1.427 & -1.000 & 1.000 & 1.000 & 0.712 \\
resampled\_stratified & gross & Beta/delta-neutral + VIX & 1.414 & 0.658 & 1.092 & 1.407 & 1.709 & 2.159 & -0.546 & 0.001 & 0.644 & 0.000 \\
resampled\_stratified & gross & Cost-aware Sortino + VIX & 1.014 & 0.539 & 0.837 & 1.022 & 1.196 & 1.454 & -1.000 & 0.001 & 1.000 & 0.958 \\
resampled\_stratified & gross & Delta-matched equities & 1.176 & 0.810 & 1.023 & 1.196 & 1.376 & 1.617 & -0.832 & 0.000 & 0.991 & 0.000 \\
resampled\_stratified & gross & Equal premium & -0.030 & -0.724 & -0.296 & -0.041 & 0.190 & 0.477 & -1.000 & 0.544 & 1.000 & 0.000 \\
resampled\_stratified & gross & Equal risk & 0.082 & -0.567 & -0.197 & 0.072 & 0.311 & 0.624 & -0.985 & 0.432 & 1.000 & 0.000 \\
resampled\_stratified & gross & Equity-option Greek Markowitz & 0.842 & 0.290 & 0.638 & 0.853 & 1.051 & 1.358 & -0.789 & 0.006 & 0.983 & 0.000 \\
resampled\_stratified & gross & Greek Markowitz + VIX & 1.374 & 0.696 & 1.086 & 1.373 & 1.657 & 2.103 & -0.571 & 0.000 & 0.701 & 0.000 \\
resampled\_stratified & gross & Underlying Markowitz & 0.909 & 0.472 & 0.748 & 0.929 & 1.089 & 1.325 & -0.260 & 0.001 & 0.000 & 0.000 \\
resampled\_stratified & gross & VIX hedge sleeve & -6.429 & -9.490 & -7.568 & -6.512 & -5.739 & -4.887 & -1.000 & 1.000 & 1.000 & 1.000 \\
\bottomrule
\end{tabular}
}}
  {\begin{tabular}{llll}\toprule Family & Basis & Strategy & Path P50 \\\midrule Missing & Run \texttt{make robustness} & -- & -- \\\bottomrule\end{tabular}}
\caption{Fixed-weight resampled-universe diagnostics.}
\label{tab:mc_resampled_universes}
\end{table}

\begin{table}[H]
\centering
\scriptsize
\IfFileExists{tables/mc_refit_stability.tex}
  {\resizebox{\textwidth}{!}{\begin{tabular}{lrrrrrrrrr}
\toprule
Strategy & Paths & OK Paths & Realized Sharpe & P05 Sharpe & P50 Sharpe & P95 Sharpe & P50 Max Drawdown & Median Gross NAV & Error Paths \\
\midrule
Greek Markowitz + VIX & 200 & 200 & 1.374 & 0.999 & 1.352 & 2.085 & -0.638 & 1.000 & 0 \\
\bottomrule
\end{tabular}
}}
  {\begin{tabular}{llll}\toprule Strategy & Paths & P50 Sharpe & Status \\\midrule Missing & Run \texttt{make robustness} & -- & -- \\\bottomrule\end{tabular}}
\caption{Refit-path stability diagnostic for Greek Markowitz + VIX.}
\label{tab:mc_refit_stability}
\end{table}

\clearpage

\subsection{Repriced synthetic universes}

The second Monte Carlo tier builds a synthetic state world. Each spot-return, implied-vol
level, and VIX-level dimension is fit with a GARCH(1,1) marginal when enough train-window
observations are available; standardized innovations are then sampled jointly by circular
block bootstrap. A Gaussian-copula sensitivity with 250 paths rank-maps innovations
through the empirical marginal distributions. The default repriced stress tier uses 1,000
paths and the same fixed train-estimated strategy weights as the headline portfolio.

Synthetic contracts are one-step, one-month options: entry premiums and payoff horizons
match. Equity options are priced with Black--Scholes at the simulated equity and IV state.
VIX options are priced with Black--76 off the simulated VX-front state at entry and settled
against the simulated VIX level at the next step. The generator's initial IV and VIX state
levels use within-train forward/backward filling only (\texttt{loc[:TRAIN\_END]}), so the
initial state never crosses into the test window.

This repriced universe is a stress world, not a replay of market option premia. Entry
premiums are model-priced from the simulated BS/Black--76 state, so the empirical
variance-risk-premium wedge embedded in observed market premia is largely absent by
construction. Strategies that harvest that premium should not be expected to retain their
realized edge in this world. Rank changes versus the resampled universes are therefore
evidence about premium-dependence, not evidence that the headline backtest failed.

\begin{table}[H]
\centering
\scriptsize
\IfFileExists{tables/mc_repriced_universes.tex}
  {\resizebox{\textwidth}{!}{\begin{tabular}{llrrrrrrr}
\toprule
Strategy & Method & Realized Sharpe & P05 Sharpe & P50 Sharpe & P95 Sharpe & P50 Max Drawdown & P Sharpe Less Than 0 & P Default \\
\midrule
Beta/delta-neutral + VIX & joint\_garch\_block & 1.414 & -3.748 & -1.073 & 0.573 & -1.000 & 0.879 & 0.741 \\
Cost-aware Sortino + VIX & joint\_garch\_block & 1.014 & -1.712 & 0.829 & 1.949 & -1.000 & 0.256 & 0.805 \\
Equal premium & joint\_garch\_block & -0.030 & -4.026 & 0.689 & 1.794 & -1.000 & 0.357 & 0.000 \\
Equal risk & joint\_garch\_block & 0.082 & -5.142 & 0.481 & 1.760 & -1.000 & 0.428 & 0.000 \\
Equity-option Greek Markowitz & joint\_garch\_block & 0.842 & -2.668 & -0.541 & 1.095 & -1.000 & 0.713 & 0.449 \\
Greek Markowitz + VIX & joint\_garch\_block & 1.374 & -3.752 & -0.919 & 0.809 & -1.000 & 0.812 & 0.681 \\
VIX hedge sleeve & joint\_garch\_block & -6.429 & -4.691 & -0.022 & 1.312 & -1.000 & 0.508 & 1.000 \\
\bottomrule
\end{tabular}
}}
  {\begin{tabular}{llll}\toprule Strategy & Method & P50 Sharpe & P Default \\\midrule Missing & Run \texttt{make robustness} & -- & -- \\\bottomrule\end{tabular}}
\caption{Repriced synthetic-universe diagnostics under the joint GARCH/block-bootstrap
state model.}
\label{tab:mc_repriced_universes}
\end{table}

\clearpage

\begin{table}[H]
\centering
\scriptsize
\IfFileExists{tables/mc_universe_comparison.tex}
  {\resizebox{\textwidth}{!}{\begin{tabular}{lrrrrrrr}
\toprule
Strategy & Sharpe Realized & Resampled P05 & Resampled P50 & Resampled P95 & Repriced P05 & Repriced P50 & Repriced P95 \\
\midrule
Beta/delta-neutral + VIX & 1.414 & 0.737 & 1.428 & 2.159 & -3.748 & -1.073 & 0.573 \\
Cost-aware Sortino + VIX & 1.014 & 0.431 & 1.006 & 1.574 & -1.712 & 0.829 & 1.949 \\
Delta-matched equities & 1.176 & 0.418 & 1.215 & 2.183 &  &  &  \\
Equal premium & -0.030 & -0.693 & -0.033 & 0.471 & -4.026 & 0.689 & 1.794 \\
Equal risk & 0.082 & -0.602 & 0.085 & 0.669 & -5.142 & 0.481 & 1.760 \\
Equity-option Greek Markowitz & 0.842 & 0.085 & 0.864 & 1.657 & -2.668 & -0.541 & 1.095 \\
Greek Markowitz + VIX & 1.374 & 0.643 & 1.391 & 2.236 & -3.752 & -0.919 & 0.809 \\
Underlying Markowitz & 0.909 & 0.133 & 0.939 & 1.763 &  &  &  \\
VIX hedge sleeve & -6.429 & -9.234 & -6.400 & -4.993 & -4.691 & -0.022 & 1.312 \\
\bottomrule
\end{tabular}
}}
  {\begin{tabular}{llll}\toprule Strategy & Realized & Resampled P50 & Repriced P50 \\\midrule Missing & Run \texttt{make robustness} & -- & -- \\\bottomrule\end{tabular}}
\caption{Comparison of realized, resampled, and repriced-universe Sharpe distributions.}
\label{tab:mc_universe_comparison}
\end{table}

\begin{figure}[H]
\centering
\IfFileExists{figures/mc_universe_sharpe_distributions.pdf}
  {\includegraphics[width=0.86\textwidth]{mc_universe_sharpe_distributions.pdf}}
  {\fbox{\parbox{0.78\textwidth}{Run \texttt{python -m research.papers.option_only_markowitz.analysis.run_empirics --stage robustness-figures} to generate the Monte Carlo universe distribution figure.}}}
\caption{Monte Carlo Sharpe distributions across resampled, regime-stratified,
repriced no-VRP stress, and refit-stability universes, with realized gross Sharpes
marked where applicable.}
\label{fig:mc_universe_sharpe_distributions}
\end{figure}

\clearpage

\subsection{Limitations}

These diagnostics broaden the evidence surface but do not remove the central research
limitations. CPCV is intentionally non-PIT. The eligible contract universe is fixed at the
headline training anchor. Resampled universes preserve observed return rows rather than
creating new contracts. Refit resampling uses a slot-calendar approximation for selected
constraint inputs. The repriced universe simplifies contract geometry and removes much of
the realized option-market premium wedge. For that reason, the section should be read as a
robustness and failure-mode map around the headline point-in-time backtest, not as a
stronger trading claim.
